\section{Evaluation and Empirical Results}
\label{sec:eval}

\subsection{Multi-Region Fault-Injection Testbed}
Each experimental seed generates 80 chronological releases across our 12-service, 6-region topology (paired identically across all evaluated strategies): 20 clean releases accompanied by exogenous regional load shocks, 15 HTTP-200 semantic payload bugs, 20 regional config/feature-flag drift faults (invisible to single-region staging), 15 production-scale tail-latency regressions, and 10 downstream brownouts. All \TotalSealedRuns{} suites (\TotalSealedRecords{} runs, \TotalSimReleases{} releases) are sealed with git commits and SHA-256 harness/artifact manifests. Disjoint 20-seed A/A calibration yields \AANoiseDetPP{}~pp detection noise and \AANoiseCyclePct{}\% cycle-time noise.

\begin{table*}[t]
\centering
\caption{Main Benchmark Across 80 Multi-Region Release Scenarios (12 Services, 6 Regions, $85\%$ Blast-Graph Recall, 25 Paired Seeds; Mean [95\% Bootstrap CI]).}
\label{tab:acv_main}
\begin{tabular}{lcccccc}
\toprule
\textbf{Strategy} & \textbf{Fault Det. (\%)} $\uparrow$ & \textbf{Sem. Rec. (\%)} $\uparrow$ & \textbf{Drift Rec. (\%)} $\uparrow$ & \textbf{False Roll. (\%)} $\downarrow$ & \textbf{Cycle (min)} $\downarrow$ & \textbf{Escapes} $\downarrow$ \\
\midrule
FixedStaging & 18.3 [18.3, 18.3] & 73.3 & 0.0 & 0.00 & 47.86 & 49.00 \\
StaticCanary5xxP99 & 3.4 [2.6, 4.3] & 0.0 & 1.2 & 24.00 & 12.61 & 42.76 \\
ShadowDiffOnly & 48.3 [46.9, 49.8] & 85.1 & 81.2 & 52.20 & 8.97 & 11.68 \\
CanaryNoCausal & 72.6 [70.6, 74.7] & 86.7 & 86.6 & 60.00 & 7.83 & 1.84 \\
\textbf{AgenticCanary} & \textbf{92.1 [90.9, 93.3]} & \textbf{93.1} & \textbf{93.8} & \textbf{9.40} & \textbf{8.21} & \textbf{3.56} \\
\bottomrule
\end{tabular}
\end{table*}


\subsection{RQ1: Fault Detection, False Rollbacks, and Release Velocity}
Table~\ref{tab:acv_main} summarizes primary results at $85\%$ blast-graph recall across 25 paired seeds:
\begin{itemize}
    \item \textbf{Vs.\ Fixed Staging}: \texttt{FixedStaging} misses \StagingDriftPct{}\% of regional config drift (0\% recall), detects only \StagingDetPct{}\% of total faults (\StagingEscapes{} escapes per 80 releases), and suffers a \StagingCycleMin{}-minute release cycle due to staging queue contention. \textbf{AgenticCanary} achieves \ACDetPct{}\% fault detection ($p\HolmPDetVsStaging{}$), \ACDriftPct{}\% config-drift recall, \ACEscapes{} escapes, and an \ACCycleMin{}-minute cycle (\CycleSpeedupVsStaging{}$\times$ faster).
    \item \textbf{Vs.\ Static 5xx/P99 Canaries}: \texttt{StaticCanary5xxP99} catches \StaticCanarySemPct{}\% of HTTP-200 semantic bugs and only \StaticCanaryDetPct{}\% of faults overall (\StaticCanaryEscapes{} escapes), while suffering $24.0\%$ false rollbacks under traffic shocks. AgenticCanary improves detection by nearly $27\times$ ($p\HolmPDetVsStatic{}$) while slashing escapes and false rollbacks.
    \item \textbf{Vs.\ Non-Causal \& Passive Shadow Baselines}: \texttt{ShadowDiffOnly} (\ShadowDetPct{}\% detection) and \texttt{CanaryNoCausal} (\NoCausalFalseRollPct{}\% false rollback) suffer severe false rollbacks during regional load spikes. AgenticCanary's DiD attribution cuts false rollbacks to \ACFalseRollPct{}\% ($p\HolmPFRVsNoCausal{}$) while achieving \ACMTTD{}~min MTTD, \ACMTTR{}~min MTTR, and blocking an average of \ACShiftLeft{} recurring faults pre-deploy via shift-left rules.
\end{itemize}

\begin{table}[t]
\centering
\caption{Component Ablation of AgenticCanary ($85\%$ Graph Recall, 25 Seeds).}
\label{tab:acv_ablation}
\begin{tabular}{lccccc}
\toprule
\textbf{Variant} & \textbf{Det\%} & \textbf{Sem\%} & \textbf{FR\%} & \textbf{Cycle} & \textbf{Esc.} \\
\midrule
\textbf{AgenticCanary} & \textbf{92.1} & \textbf{93.1} & \textbf{9.4} & \textbf{8.21} & \textbf{3.56} \\
No-NoBlastGraph & 95.4 & 99.7 & 8.2 & 9.26 & 1.60 \\
No-NoShadowProbe & 75.0 & 24.5 & 9.2 & 9.47 & 12.88 \\
No-NoCausalDiD & 76.7 & 88.8 & 60.0 & 6.95 & 1.76 \\
No-NoShiftLeft & 90.2 & 91.5 & 9.4 & 9.28 & 4.32 \\
\bottomrule
\end{tabular}
\end{table}


\subsection{RQ2--RQ3: Ablation, Graph Recall, and Traffic Shock Stress}
Table~\ref{tab:acv_ablation} confirms every mechanism contributes materially: removing shadow semantic probing drops semantic recall from \ACSemPct{}\% to \AblNoShadowSemPct{}\%; removing DiD causal attribution spikes false rollbacks from \ACFalseRollPct{}\% to \AblNoCausalFRPct{}\%; removing shift-left memory increases escaped outages from \ACEscapes{} to \AblNoShiftEscapes{}; and removing blast-graph targeting slows cycle time from \ACCycleMin{} to \AblNoBlastCycle{}~min.

Table~\ref{tab:acv_recall} validates non-circularity under incomplete blast graphs: even at $50\%$ graph recall (half of fault cones hidden from pre-deploy analysis), AgenticCanary sustains \RecallFiftyACDetPct{}\% detection vs.\ \RecallFiftyShadowDetPct{}\% for passive shadow diffing. Under $2\times$ extreme exogenous traffic shocks (\texttt{exp4\_shock\_stress\_v1}), AgenticCanary holds false rollbacks to \ShockTwoACFRPct{}\% compared to \ShockTwoNoCausalFRPct{}\% for \texttt{CanaryNoCausal}.

\begin{table}[t]
\centering
\caption{Non-Circularity Test: Fault Detection (\%) vs.\ Blast-Graph Recall ($50\%$--$100\%$, 20 Seeds).}
\label{tab:acv_recall}
\begin{tabular}{lcccc}
\toprule
\textbf{Graph Recall} & \textbf{Staging} & \textbf{StaticCan.} & \textbf{Shadow} & \textbf{AgenticCan.} \\
\midrule
50\% & 18.3 & 3.2 & 48.3 & \textbf{79.1} \\
70\% & 18.3 & 3.2 & 48.3 & \textbf{86.9} \\
85\% & 18.3 & 3.2 & 48.3 & \textbf{92.4} \\
100\% & 18.3 & 3.2 & 48.3 & \textbf{98.8} \\
\bottomrule
\end{tabular}
\end{table}

